%! Introduction to Eigenvalues — Unit 6, Lesson 6.1 — Homework
\documentclass[10pt]{article}
\usepackage{linalg-article}
\usepackage{linalg-boxes}
\newcommand{\TallMath}[1]{$\displaystyle #1\rule[-1.4em]{0pt}{3.2em}$}
\newcommand{\xx}{\mathbf{x}}
\newcommand{\zero}{\mathbf{0}}

\begin{document}

\pageheader{Unit 6, Lesson 6.1}{Homework}
\namedateperiod

\vspace{0.12in}

\begin{notesbox}{Practice}
The method is two steps. \textbf{Step 1:} solve $\det(A-\lambda I)=0$ for the \textbf{eigenvalues}.
\textbf{Step 2:} for each $\lambda$, solve $(A-\lambda I)\xx=\zero$ for an \textbf{eigenvector}. Check
with the trace ($=\lambda_1+\lambda_2$) and determinant ($=\lambda_1\lambda_2$).

\begin{enumerate}[leftmargin=1.9em, itemsep=12pt]
  \item \textbf{Full method.} For $A=\begin{bmatrix} 5 & 2 \\ 2 & 5 \end{bmatrix}$, find both
        eigenvalues, then an eigenvector for each. \writelines{3}
  \item \textbf{Read the eigenvalues off.} State the eigenvalues of each matrix (no full work needed ---
        look at the diagonal), and say why.
    \begin{enumerate}[label=(\alph*), leftmargin=1.8em, itemsep=4pt]
      \item $A=\begin{bmatrix} 6 & 0 \\ 0 & 2 \end{bmatrix}$ \hfill
      \item $A=\begin{bmatrix} 4 & 5 \\ 0 & 1 \end{bmatrix}$ \;(triangular)
    \end{enumerate}
    \vspace{1.0cm}
  \item \textbf{A negative eigenvalue.} For $A=\begin{bmatrix} 1 & 2 \\ 2 & 1 \end{bmatrix}$, find both
        eigenvalues and an eigenvector for each. What does the \emph{negative} eigenvalue do to its
        direction? \writelines{3}
  \item \textbf{Trace \& determinant check.} For $A=\begin{bmatrix} 4 & 1 \\ 2 & 3 \end{bmatrix}$ (the
        warm-up matrix), find both eigenvalues. Then verify that their \emph{sum} equals the trace and
        their \emph{product} equals $\det A$. \writelines{2}
  \item \textbf{Explain.} Why does solving $\det(A-\lambda I)=0$ give the eigenvalues, and why must an
        eigenvector be \emph{nonzero}? \writelines{2}
\end{enumerate}
\end{notesbox}

\begin{extensionbox}
\textbf{Eigenvalues of a $3\times3$ triangular matrix.} The method scales up: for a triangular matrix,
$\det(A-\lambda I)$ is the product down the diagonal. Find the eigenvalues of
\[
  A=\begin{bmatrix} 2 & 1 & 7 \\ 0 & 3 & 4 \\ 0 & 0 & 5 \end{bmatrix}.
\]
(\emph{Hint:} $\det(A-\lambda I)=(2-\lambda)(3-\lambda)(5-\lambda)$.) \writelines{2}
\end{extensionbox}

\begin{spiralbox}
\textbf{Coming next --- Lesson 6.2: Diagonalizing a Matrix.} You can now \emph{find} a matrix's
eigenvalues and eigenvectors. Next we line up the eigenvectors into a matrix and use them to
\emph{rebuild} $A$ as $A=X\Lambda X^{-1}$ --- turning powers, and hard matrix problems, into easy work
on the numbers $\lambda$.
\end{spiralbox}

\end{document}
